% Notes of the Montreal seminar, 1970 (cote 129, typed part): archive cote 129, site page 17 (PDF page 18).
% typescript. Draft compared: none.
% AI-transcribed by Claude Opus 5.5 (claude-opus-5-5, adaptive thinking, effort high), Message Batches API, batch mtl129_full (msgbatch_01AGnDKHwy6ATCM7ehUa6263), 2026-11-13; not yet checked by eye.
\scanpage{129/17}
Les morphismes de restriction
\[
R_n : \mathbb{W}_{n+1} \to \mathbb{W}_n \qquad (x_1,\ldots,x_n,x_{n+1}) \to (x_1,\ldots,x_n)
\]
sont alors des homomorphismes d'anneaux. En effet, il résulte de l'uni-%
cité de la structure d'anneau sur $\mathbb{W}_n$ qu'un morphisme $f : X \to \mathbb{W}_n$ est
un morphisme d'anneaux (resp. de groupes) si et seulement si $\phi_i \circ f : X \to \underline{O}$
est un homomorphisme d'anneaux (resp. de groupes pour tout $i$, $1 \leq i \leq n$ .
Ici, on a $\phi_i R_n = \phi_i$ , qui est par définition un homomorphisme d'anneaux.

On définit le schéma en anneau des vecteurs de Witt $\mathbb{W}$ par
\[
\mathbb{W} = \lim_{\substack{\leftarrow \\ n}} (\mathbb{W}_n , R_n) \quad .
\]
Le schéma sous-jacent à $\mathbb{W}$ est donc en particulier isomorphe à $\mathbb{E}^{\mathbb{N}}$ .

2. \uline{Quelques morphismes remarquables}

2.1. On définit
\[
V = V_n : \mathbb{W}_n \to \mathbb{W}_n \qquad (x_1,x_2,\ldots,x_{n-1},x_n) \to (0,x_1,\ldots,x_{n-1})
\]
\[
T = T_n : \mathbb{W}_n \to \mathbb{W}_{n+1} \qquad (x_1,x_2,\ldots,x_n) \to (0,x_1,\ldots,x_n) \ .
\]

Le morphisme $T$ est \uline{additif} ; on a en effet $\phi_i\, T = p\, \phi_{i-1}$
pour $1 \leq i \leq n+1$ .

On en déduit que $V_n = R_n\, T_n = T_{n-1}\, R_{n-1}$ est aussi un morphis-%
me \uline{additif}.

\mspage{\hand{17}}
